% !TEX program = xelatex
% Lernplatt - by Can Siebert
% Kurs 71 (Dig 42) - Tag 3 - Lehrskript
% Vom Prozessbild zur Steuerung: Stakeholder - KPIs - SOA - Workflow
%
% Self-contained xelatex source. Kein biblatex, keine externen Bilder.

\documentclass[11pt,a4paper,oneside]{article}

% --- Encoding & Sprache --------------------------------------------------
\usepackage{polyglossia}
\setdefaultlanguage{german}

% --- Geometrie -----------------------------------------------------------
\usepackage[a4paper,top=2cm,bottom=2cm,outer=2.5cm,inner=2.5cm,bindingoffset=1.5cm]{geometry}

% --- Fonts ---------------------------------------------------------------
\usepackage{fontspec}
\defaultfontfeatures{Ligatures=TeX}

\IfFontExistsTF{Charter}{%
  \setmainfont{Charter}[Numbers=OldStyle]%
}{%
  \IfFontExistsTF{Noto Serif}{%
    \setmainfont{Noto Serif}%
  }{%
    \setmainfont{Liberation Serif}%
  }%
}

\IfFontExistsTF{Source Sans 3}{%
  \setsansfont{Source Sans 3}%
}{%
  \IfFontExistsTF{Source Sans Pro}{%
    \setsansfont{Source Sans Pro}%
  }{%
    \setsansfont{Liberation Sans}%
  }%
}

\newcommand{\caveatfontfallback}{\itshape}
\IfFontExistsTF{Caveat}{%
  \newfontfamily\caveatfont{Caveat}[Scale=1.15]%
}{%
  \newcommand\caveatfont{\caveatfontfallback}%
}

\setmonofont{Liberation Mono}

% --- Mikro-Typografie ----------------------------------------------------
\usepackage{microtype}
\linespread{1.15}

% --- Farben & Boxen ------------------------------------------------------
\usepackage{xcolor}
\definecolor{lpInk}{RGB}{20,28,38}
\definecolor{kurs71}{HTML}{9D4A1A}
\definecolor{lpAccent}{HTML}{9D4A1A}
\definecolor{lpAccentLight}{RGB}{248,232,218}
\definecolor{lpAmber}{RGB}{222,138,30}
\definecolor{lpAmberLight}{RGB}{255,243,217}
\definecolor{lpAmberBorder}{RGB}{196,118,18}
\definecolor{lpGray}{RGB}{96,104,114}
\definecolor{lpGrayLight}{RGB}{238,240,243}

\usepackage[most]{tcolorbox}
\tcbuselibrary{breakable,skins}

\newtcolorbox{eselsbox}[1][Eselsbrücke]{%
  enhanced, breakable,
  colback=lpAmberLight,
  colframe=lpAmberBorder,
  boxrule=0.6pt,
  arc=2pt,
  borderline={0.4pt}{0pt}{lpAmberBorder, dashed},
  left=10pt,right=10pt,top=8pt,bottom=8pt,
  fonttitle=\sffamily\bfseries\color{lpAmberBorder},
  coltitle=lpAmberBorder,
  title={#1},
  attach boxed title to top left={xshift=8pt,yshift=-8pt},
  boxed title style={size=small, colback=lpAmberLight, colframe=lpAmberBorder, boxrule=0.4pt, arc=1pt},
  top=14pt
}

\newtcolorbox{merkbox}[1][Merksatz]{%
  enhanced, breakable,
  colback=lpAccentLight,
  colframe=lpAccent,
  boxrule=0.6pt,
  arc=2pt,
  left=10pt,right=10pt,top=8pt,bottom=8pt,
  fonttitle=\sffamily\bfseries\color{white},
  coltitle=white,
  title={#1},
  attach boxed title to top left={xshift=8pt,yshift=-8pt},
  boxed title style={size=small, colback=lpAccent, colframe=lpAccent, boxrule=0.4pt, arc=1pt},
  top=14pt
}

% --- Listen --------------------------------------------------------------
\usepackage{enumitem}
\setlist{itemsep=2pt, topsep=4pt, parsep=0pt}
\setlist[itemize,1]{label={\textbullet},leftmargin=1.6em}
\setlist[itemize,2]{label={\textendash},leftmargin=1.4em}
\setlist[enumerate,1]{label=\arabic*., leftmargin=1.8em}

% --- Tabellen ------------------------------------------------------------
\usepackage{tabularx}
\usepackage{array}
\usepackage{booktabs}
\usepackage{longtable}
\usepackage{multirow}

% --- Euro-Zeichen --------------------------------------------------------
\usepackage{eurosym}

% --- Kopf-/Fusszeilen ----------------------------------------------------
\usepackage{fancyhdr}
\usepackage{lastpage}
\pagestyle{fancy}
\fancyhf{}
\renewcommand{\headrulewidth}{0.3pt}
\renewcommand{\footrulewidth}{0pt}
\fancyhead[L]{\footnotesize\sffamily\color{lpGray}Lernplatt \textperiodcentered{} Kurs 71 \textperiodcentered{} Tag 3}
\fancyhead[R]{\footnotesize\sffamily\color{lpGray}Stakeholder \textperiodcentered{} KPIs \textperiodcentered{} SOA \textperiodcentered{} Workflow}
\fancyfoot[L]{\footnotesize\sffamily\color{lpGray}\textcopyright{} Can Siebert}
\fancyfoot[R]{\footnotesize\sffamily\color{lpGray}\thepage{} / \pageref{LastPage}}

\fancypagestyle{plain}{%
  \fancyhf{}%
  \renewcommand{\headrulewidth}{0pt}%
  \fancyfoot[C]{\footnotesize\sffamily\color{lpGray}\thepage{} / \pageref{LastPage}}%
}

% --- Section-Styling -----------------------------------------------------
\usepackage[explicit]{titlesec}

\titleformat{\section}[block]
  {\sffamily\Large\bfseries\color{lpAccent}}
  {\thesection.}{0.6em}{#1}
  [{\vspace{2pt}\color{lpAccent}\titlerule[0.6pt]}]

\titleformat{\subsection}[block]
  {\sffamily\large\bfseries\color{lpInk}}
  {\thesubsection}{0.6em}{#1}

\titleformat{\subsubsection}[block]
  {\sffamily\normalsize\bfseries\color{lpInk}}
  {}{0pt}{#1}

\titlespacing*{\section}{0pt}{14pt}{8pt}
\titlespacing*{\subsection}{0pt}{10pt}{4pt}
\titlespacing*{\subsubsection}{0pt}{8pt}{2pt}

% --- Hyperlinks ----------------------------------------------------------
\usepackage[hidelinks,unicode]{hyperref}

% --- TikZ ----------------------------------------------------------------
\usepackage{tikz}
\usetikzlibrary{shapes.geometric,arrows.meta,positioning,calc,fit,backgrounds}
\definecolor{lpGreenLight}{RGB}{225,238,225}

% --- TOC -----------------------------------------------------------------
\renewcommand{\contentsname}{\sffamily\Large\bfseries\color{lpAccent}Inhalt}

% --- Helfer --------------------------------------------------------------
\newcommand{\akzent}[1]{{\caveatfont\color{lpAmber}#1}}
\newcommand{\fachbegriff}[1]{\textbf{#1}}
\newcommand{\quelle}[1]{{\footnotesize\color{lpGray}(#1)}}

% =========================================================================
\begin{document}

% --- COVER ---------------------------------------------------------------
\begin{titlepage}
  \pagestyle{empty}
  \vspace*{1.5cm}
  \noindent{\sffamily\footnotesize\color{lpGray}LERNPLATT \textperiodcentered{} BY CAN SIEBERT}\\[2pt]
  \noindent{\color{lpAccent}\rule{\linewidth}{1.2pt}}\\[4pt]
  \noindent{\sffamily\footnotesize\color{lpGray}MODUL 2 \textperiodcentered{} TAG 3 VON 4 \textperiodcentered{} KURS 71 (Dig 42) \textperiodcentered{} 1.\,Juni 2026}

  \vspace{3cm}
  {\sffamily\fontsize{30}{34}\selectfont\bfseries\color{lpInk}Vom Prozessbild\\zur Steuerung\par}

  \vspace{0.7cm}
  {\sffamily\large\color{lpGray}Stakeholder, KPIs, SOA-Services und Workflow-Orchestrierung --\\die vier Hebel eines steuerbaren Prozess-Operating-Models.\par}

  \vspace{1.5cm}
  {\caveatfont\LARGE\color{lpAmber}Lehrskript -- Tag 3 von 4}\par

  \vfill

  \noindent\begin{minipage}{0.55\linewidth}
    {\sffamily\footnotesize\color{lpGray}Lernpfad:}\\
    {\sffamily\small\color{lpInk}Wissen \textrightarrow{} Anwendung \textrightarrow{} Management}\\[10pt]
    {\sffamily\footnotesize\color{lpGray}Bearbeitungsdauer:}\\
    {\sffamily\small\color{lpInk}1 Kurstag}
  \end{minipage}\hfill
  \begin{minipage}{0.4\linewidth}
    \raggedleft
    {\sffamily\footnotesize\color{lpGray}Dozent:}\\
    {\sffamily\small\color{lpInk}Can Siebert}\\[10pt]
    {\sffamily\footnotesize\color{lpGray}Format:}\\
    {\sffamily\small\color{lpInk}Theorie \textperiodcentered{} Fallstudie \textperiodcentered{} Coaching}
  \end{minipage}

  \vspace{0.5cm}
  \noindent{\color{lpAccent}\rule{\linewidth}{0.6pt}}
\end{titlepage}

% --- LERNZIELE -----------------------------------------------------------
\section*{Lernziele dieses Tages}
\addcontentsline{toc}{section}{Lernziele dieses Tages}
\thispagestyle{plain}

\noindent An den ersten beiden Tagen haben Sie Prozesse \emph{modelliert} (BPMN/EPC, UML, VSM, Process Mining). Heute geht es um die Frage: \emph{Wie wird aus einem Modell ein Steuerungsinstrument?} Vier Hebel: Stakeholder-Klärung, KPI-Architektur, SOA-Servicegrenzen und Workflow-Orchestrierung. Ohne diese vier bleibt jedes Diagramm Papier.

\subsection*{L1 -- Wissen (Reproduktion und Einordnung)}
\begin{itemize}
  \item Stakeholder-Analyse strukturieren: 2$\times$2-Matrix (Einfluss $\times$ Interesse), Einbindungsstufen.
  \item KPI-Frameworks unterscheiden: Outcome, Output, Prozess, Compliance/Risiko -- Leading vs.\ Lagging.
  \item Goodhart-Effekt und KPI-Steckbrief als Schutz vor Fehlsteuerung kennen.
  \item SOA-Grundbegriffe: Service, Kohäsion, Kopplung, Orchestrierung vs.\ Choreographie.
  \item Workflow-Bausteine: Aufgaben, Regeln, Rollen, Eskalation, SLA, Audit Trail.
\end{itemize}

\subsection*{L2 -- Anwendung (Transfer auf konkrete Situationen)}
\begin{itemize}
  \item Eine Stakeholder-Map mit Einbindungsstrategie für ein konkretes Vorhaben erstellen.
  \item Ein KPI-Set für einen Prozess entwerfen -- inklusive vollständiger KPI-Steckbriefe.
  \item Aus einem Prozess Service-Kandidaten ableiten (Input/Output/Owner).
  \item Eine Workflow-Skizze mit SLAs und Eskalationen entwerfen.
\end{itemize}

\subsection*{L3 -- Management (Entscheidung und Verantwortung)}
\begin{itemize}
  \item Eine Steuerungsarchitektur konzipieren, die Entscheidungen tatsächlich stützt -- nicht nur Reports erzeugt.
  \item Zielkonflikte (Speed vs.\ Quality, Standardisierung vs.\ Flexibilität, Transparenz vs.\ Messangst) transparent moderieren.
  \item KPI-Manipulationsrisiken vorab adressieren (Goodhart-Schutz durch Kopplungen).
  \item Servicegrenzen strategisch entscheiden -- Over-Engineering vs.\ unklare Verantwortung.
\end{itemize}

\begin{merkbox}[Roter Faden]
Ein Prozessmodell ohne klare Stakeholder, ohne wirksame KPIs, ohne saubere Servicegrenzen und ohne Workflow-Steuerung ist ein \emph{Diagramm}. Mit diesen vier Hebeln wird es ein \emph{Operating Model}. Senior-Verantwortung heißt, alle vier Hebel zu beherrschen -- und ihre Wechselwirkungen.
\end{merkbox}

\clearpage

% --- 1. EINSTIEG ---------------------------------------------------------
\section{Einstieg: Vom Diagramm zur Steuerung}

Bisher haben wir Prozesse beschrieben (BPMN, EPC, UML), Wertströme analysiert (VSM) und Datenrealität geprüft (Process Mining). Das sind Werkzeuge zum \emph{Verstehen}. Heute geht es um Werkzeuge zum \emph{Steuern}. Der Unterschied ist groß: ein verstandener Prozess kann immer noch chaotisch laufen. Ein gesteuerter Prozess hat klare Entscheidungswege, messbare Wirkung und nachweisbare Verantwortung.

\subsection{Was ``Steuerung'' tatsächlich bedeutet}

Steuerung ist mehr als ein Dashboard. Sie umfasst vier Elemente, die zusammenwirken müssen: \quelle{vgl.\ Dumas et al., 2018, Kap.~6--7}

\begin{enumerate}
  \item \fachbegriff{Stakeholder-Klärung.} Wer ist betroffen, wer entscheidet, wer wird informiert, wer kann stoppen?
  \item \fachbegriff{KPI-Architektur.} Welche Kennzahlen geben Auskunft über die richtige Frage -- und wie verhindern wir, dass die Kennzahl selbst zum Problem wird (Goodhart-Effekt)?
  \item \fachbegriff{SOA-Servicegrenzen.} Welche Fähigkeiten kapseln wir zu Services, die wiederverwendbar und klar abgegrenzt sind?
  \item \fachbegriff{Workflow-Orchestrierung.} Welches System steuert die Tasks, Eskalationen und Übergaben aktiv -- statt sie dem Zufall zu überlassen?
\end{enumerate}

Fehlt einer dieser vier Hebel, bleibt Steuerung Behauptung: ``Wir haben Prozesse'' (ohne klare Owner), ``Wir messen alles'' (mit Vanity Metrics), ``Wir nutzen Services'' (ohne Kapselung), ``Wir haben einen Workflow'' (in PowerPoint, nicht in einem System).

\subsection{Vier typische Pathologien ohne Steuerung}

\begin{itemize}
  \item \fachbegriff{Schatten-Entscheidungen.} Jeder entscheidet ad hoc, ohne klare Mandate. Niemand kann später nachvollziehen, wer warum entschieden hat.
  \item \fachbegriff{KPI-Aktivismus.} Es werden 30 Kennzahlen gemessen, aber keine führt zu einer konkreten Steuerungs-Reaktion.
  \item \fachbegriff{Spaghetti-Integration.} Jedes IT-System spricht direkt mit jedem anderen -- ohne Abstraktion. Änderungen werden teuer und riskant.
  \item \fachbegriff{Eskalation per Zufall.} Wann eskaliert wird, hängt davon ab, wer gerade besonders laut ist -- nicht von einer klaren Regel.
\end{itemize}

\begin{eselsbox}[Eselsbrücke: \akzent{S-K-S-W}]
Die vier Hebel der Prozess-Steuerung:\\[2pt]
\textbf{S}takeholder -- wer hat Einfluss, Interesse, Mandate?\\
\textbf{K}PIs -- welche Kennzahlen stützen welche Entscheidung?\\
\textbf{S}ervices -- welche Fähigkeiten kapseln wir wie?\\
\textbf{W}orkflow -- wer steuert Tasks aktiv mit Regeln und SLAs?\\[2pt]
\emph{``S-K-S-W'' -- ohne diese vier ist ein Modell nur ein Bild.}
\end{eselsbox}

% --- 2. STAKEHOLDER-ANALYSE ----------------------------------------------
\section{Stakeholder-Analyse: Wer ist betroffen, wer entscheidet?}

\fachbegriff{Stakeholder} sind alle Personen oder Gruppen, die ein Vorhaben beeinflussen können oder davon betroffen sind. Der Begriff geht auf Freeman zurück, der ihn als systematischen Bestandteil strategischen Managements etabliert hat. \quelle{Freeman, 1984}

\subsection{Was eine Stakeholder-Analyse leisten muss}

Eine gute Stakeholder-Analyse beantwortet vier Fragen für jeden Stakeholder:

\begin{enumerate}
  \item \textbf{Wer?} -- konkret benannt: Rolle (nicht Person), Abteilung, externe Beteiligte.
  \item \textbf{Welches Interesse?} -- was will der Stakeholder erreichen oder verhindern?
  \item \textbf{Welcher Einfluss?} -- kann der Stakeholder das Vorhaben stoppen, beschleunigen, verändern?
  \item \textbf{Welche Einbindung?} -- informieren, konsultieren, mitentscheiden oder verantworten?
\end{enumerate}

\subsection{Die 2$\times$2-Matrix (Einfluss $\times$ Interesse)}

Die bekannteste Stakeholder-Matrix geht auf Mendelow zurück -- sie ordnet Stakeholder nach Einfluss (Power) und Interesse: \quelle{Mendelow, 1981}

\begin{tabularx}{\linewidth}{@{}>{\bfseries}lXX@{}}
\toprule
       & \textbf{Niedriges Interesse} & \textbf{Hohes Interesse} \\
\midrule
Niedriger Einfluss & \textbf{Minimal beachten} -- regelmäßige Information genügt & \textbf{Informiert halten} -- aktiv kommunizieren, aber keine Mitsprache nötig \\
Hoher Einfluss & \textbf{Zufrieden halten} -- Erwartungen managen, keine Reibung erzeugen & \textbf{Eng managen} -- Mitentscheidung, intensiver Austausch, gemeinsame Ziele \\
\bottomrule
\end{tabularx}

\medskip

\noindent Die häufigste Falle: \emph{alle Stakeholder gleich behandeln.} Niedrige Einfluss/Interesse-Stakeholder mit gleichem Aufwand zu beteiligen wie ``eng managen''-Stakeholder verschwendet Ressourcen. Hohe Einfluss-Stakeholder zu ignorieren erzeugt Stopper.

\subsection{Vier Einbindungsstufen}

Statt ``mehr oder weniger einbinden'' arbeiten wir mit vier klaren Stufen:

\begin{enumerate}
  \item \fachbegriff{Informieren} -- der Stakeholder erfährt von Ergebnissen. Keine Mitsprache.
  \item \fachbegriff{Konsultieren} -- der Stakeholder kann Feedback geben. Die Entscheidung bleibt beim Owner.
  \item \fachbegriff{Mitentscheiden} -- der Stakeholder ist Teil des Entscheidungs-Gremiums. Konsens oder Mehrheit.
  \item \fachbegriff{Verantworten} -- der Stakeholder ist Owner. Trägt die Entscheidung allein und kann eskaliert werden.
\end{enumerate}

Die Einbindungsstufe wird \emph{vor} dem Vorhaben festgelegt -- nicht im Konflikt. Wer hinterher entscheidet, ist nicht steuerbar.

\subsection{Typische Stakeholder-Konflikte}

In Prozess-Verbesserungs-Projekten tauchen drei Konfliktarten immer wieder auf:

\begin{itemize}
  \item \fachbegriff{Speed vs.\ Kontrolle.} F\&E will Innovationsfreiheit, QM will Gates. Vertrieb will Schnelligkeit, Compliance will Prüfung.
  \item \fachbegriff{Standardisierung vs.\ Flexibilität.} IT will einheitliche Prozesse, Fachbereiche wollen Sondererfüllungen für ``ihre'' Kunden.
  \item \fachbegriff{Transparenz vs.\ Angst vor Messung.} Management will Kennzahlen, operative Teams fürchten, dass Kennzahlen gegen sie verwendet werden.
\end{itemize}

Diese Konflikte sind \emph{nicht zu lösen} im Sinne von ``richtige Antwort''. Sie sind zu \emph{moderieren} -- sichtbar zu machen, in Trade-offs zu übersetzen, mit klarer Entscheidung zu versehen. Wer Konflikte verdrängt, verschiebt sie nur in den Untergrund, wo sie das Projekt zerlegen.

\begin{merkbox}[Stopper-Stakeholder identifizieren]
Die wichtigste Stakeholder-Frage lautet nicht ``wer hilft uns?'' sondern \emph{``wer kann uns stoppen?''}. Diese Stakeholder müssen aktiv eingebunden werden -- nicht weil sie immer Recht haben, sondern weil ihre Ablehnung das Vorhaben kostet, was wir an Zeit, Geld und Glaubwürdigkeit haben.
\end{merkbox}

% --- 3. KPI-FRAMEWORKS ---------------------------------------------------
\section{KPI-Frameworks und Performance-Messung}

Wer steuern will, muss messen. Wer falsch misst, steuert falsch. Die KPI-Disziplin trennt steuernde Plattform-Teams von berichtenden.

\subsection{Vier KPI-Ebenen}

Wirksame KPI-Architekturen unterscheiden vier Ebenen: \quelle{vgl.\ Parmenter, 2015, Kap.~1; Kaplan \& Norton, 1996}

\begin{tabularx}{\linewidth}{@{}>{\bfseries}lXX@{}}
\toprule
Ebene & Misst & Beispiel \\
\midrule
Outcome     & Welcher Kundennutzen entsteht?         & Kundenzufriedenheit nach Pilot, NPS \\
Output      & Welche Menge wird produziert?           & Anzahl freigegebener Konzepte/Quartal \\
Prozess     & Wie effizient ist der Ablauf?          & Lead Time, Rework-Rate, Kosten/Vorgang \\
Compliance  & Welche Risiken/Auflagen werden erfüllt? & Audit-Findings, Fehlerquote, Eskalationen \\
\bottomrule
\end{tabularx}

\medskip

\noindent Eine ausgewogene KPI-Architektur deckt alle vier Ebenen ab. Wer nur Output misst, optimiert für Menge auf Kosten von Qualität. Wer nur Compliance misst, bekommt Bürokratie ohne Wertbeitrag.

\subsection{Leading vs.\ Lagging Indicators}

Eine zweite wichtige Unterscheidung: \quelle{Parmenter, 2015}

\begin{itemize}
  \item \fachbegriff{Lagging Indicators} messen \emph{Ergebnisse} (was passiert ist) -- z.\,B.\ Umsatz, Lead Time, Rework. Sie sind präzise, aber spät -- man kann nicht mehr eingreifen.
  \item \fachbegriff{Leading Indicators} messen \emph{Vorboten} (was passieren wird) -- z.\,B.\ Anteil ``Ideen mit Business Case'', frühe Qualitätssignale. Sie sind ungenauer, aber rechtzeitig -- man kann noch reagieren.
\end{itemize}

Eine gute KPI-Architektur hat \emph{beide}. Lagging KPIs für Verantwortlichkeit (was haben wir erreicht?), Leading KPIs für Steuerung (wo müssen wir nachsteuern?).

\subsection{Der Goodhart-Effekt}

Charles Goodhart hat 1975 eine Regelmäßigkeit formuliert, die heute oft so paraphrasiert wird: \emph{``When a measure becomes a target, it ceases to be a good measure.''} \quelle{nach Goodhart, in Strathern, 1997}

Beispiele aus Prozess-Steuerung:

\begin{itemize}
  \item Wenn ``Anzahl freigegebener Konzepte'' das Hauptziel ist, werden mehr unreife Konzepte freigegeben.
  \item Wenn ``Lead Time'' das Hauptziel ist, werden Vorgänge frühzeitig geschlossen ohne sauberen Qualitätsnachweis.
  \item Wenn ``Audit-Findings = 0'' das Hauptziel ist, werden Probleme nicht mehr gemeldet -- nur noch verdeckt.
\end{itemize}

Schutz gegen den Goodhart-Effekt: \emph{Kopplungen.} Jede einzeln optimierbare KPI wird mit einer Gegen-KPI gekoppelt, die das Manipulationsrisiko adressiert. ``Lead Time'' wird gekoppelt mit ``Rework-Rate'' und ``Gate-Pass-Rate beim ersten Durchlauf''.

\subsection{Der KPI-Steckbrief}

Eine wirksame KPI braucht einen vollständigen Steckbrief -- nicht nur einen Namen: \quelle{Parmenter, 2015, Kap.~6}

\begin{tabularx}{\linewidth}{@{}>{\bfseries}lX@{}}
\toprule
Feld & Beschreibung \\
\midrule
Name              & Eindeutig, kurz, business-verständlich \\
Definition        & Was wird genau gemessen? Was wird ausgeschlossen? \\
Formel            & Wie berechnet sich der Wert? Quotient, Median, Summe? \\
Datenquelle       & Aus welchem System? Wer pflegt die Daten? \\
Messfrequenz      & Täglich, wöchentlich, monatlich, quartalsweise? \\
Owner             & Wer ist verantwortlich für Pflege und Reaktion? \\
Ziel / Schwelle   & Welcher Wert ist ``gut''? Was ist eine Eskalations-Schwelle? \\
Gewünschtes Verhalten & Welche Reaktion soll die KPI auslösen? \\
Manipulations\-risiko & Wie könnte die KPI verfälscht werden? \\
Gegenmaßnahme     & Welche andere KPI schützt vor Manipulation? \\
\bottomrule
\end{tabularx}

\medskip

\noindent Wer einen KPI-Steckbrief nicht ausfüllen kann, hat keine KPI -- nur eine Zahl. Der Steckbrief ist Pflicht, weil er die unbequemen Fragen erzwingt: \emph{Wer macht was, wenn die KPI eine bestimmte Schwelle reißt?}

\begin{eselsbox}[Eselsbrücke: \akzent{N-D-F-D-M-O-Z-V-M-G}]
Die zehn Felder eines KPI-Steckbriefs:\\[2pt]
\textbf{N}ame -- \textbf{D}efinition -- \textbf{F}ormel -- \textbf{D}atenquelle -- \textbf{M}essfrequenz -- \textbf{O}wner -- \textbf{Z}iel -- \textbf{V}erhalten -- \textbf{M}anipulationsrisiko -- \textbf{G}egenmaßnahme.\\[2pt]
\emph{Wer alle zehn beantworten kann, hat eine steuerbare KPI. Wer es nicht kann, hat eine schöne Zahl.}
\end{eselsbox}

% --- 4. SOA -------------------------------------------------------------
\section{SOA -- Servicegrenzen sauber ziehen}

\fachbegriff{Service-Oriented Architecture (SOA)} ist ein Architekturprinzip, das Geschäftsfähigkeiten in klar abgegrenzten, wiederverwendbaren Services kapselt. Auch wenn der Begriff aus den frühen 2000er Jahren stammt, sind seine Grundideen heute aktueller denn je -- nur unter neuen Namen (Microservices, APIs, Domain-Driven Design). \quelle{Erl, 2005}

\subsection{Was ein Service ist}

Ein \fachbegriff{Service} im SOA-Sinne ist eine klar abgegrenzte Fähigkeit:

\begin{itemize}
  \item Hat einen klaren Zweck (eine fachliche Aufgabe).
  \item Hat einen klaren Vertrag (Input/Output).
  \item Ist wiederverwendbar (verschiedene Prozesse können ihn nutzen).
  \item Hat einen klaren Owner (eine Rolle ist verantwortlich).
\end{itemize}

Beispiel: ``Kundenstammdaten prüfen'' ist ein Service. Eingang: Kunden-ID; Ausgang: Status der Datenvollständigkeit + Liste fehlender Felder; Owner: Stammdaten-Team. Dieser Service kann von verschiedenen Prozessen genutzt werden -- Bestellung, Vertrag, Reklamation -- ohne dass jeder Prozess eigene Logik dafür baut.

\subsection{Die zwei Qualitätskriterien: Kohäsion und Kopplung}

Service-Qualität misst sich an zwei Kriterien: \quelle{Erl, 2005, Kap.~5--6}

\begin{itemize}
  \item \fachbegriff{Hohe Kohäsion} -- alles in einem Service gehört fachlich zusammen. Der Service tut \emph{eine} Sache, nicht fünf.
  \item \fachbegriff{Geringe Kopplung} -- der Service ist möglichst unabhängig von anderen Services. Änderungen in einem Service brechen nicht zwangsläufig andere.
\end{itemize}

Diese beiden Prinzipien sind 50 Jahre alt (David Parnas, 1972) und gelten unverändert. Sie sind das eigentliche Geheimnis tragfähiger IT-Architekturen.

\subsection{Wann lohnt sich ein Service?}

Nicht jede Funktion verdient einen eigenen Service. Die Faustregel:

\begin{itemize}
  \item \textbf{Wiederverwendung} -- wird die Funktion von mehreren Prozessen genutzt?
  \item \textbf{Komplexität} -- ist die Logik komplex genug, dass Kapselung Wert schafft?
  \item \textbf{Compliance / Standardisierung} -- gibt es Vorgaben (rechtlich, intern), die zentrale Umsetzung verlangen?
\end{itemize}

Wenn keine dieser drei Bedingungen zutrifft, ist ein eigener Service Over-Engineering. Dann reicht eine Bibliotheksfunktion oder ein einfacher Workflow-Schritt.

\subsection{Häufige SOA-Fehler}

\begin{enumerate}
  \item \fachbegriff{Mega-Service.} Der Service tut zu viel -- ``Order Service'' mit 30 Funktionen. Kohäsion: niedrig.
  \item \fachbegriff{Mikro-Mikro-Service.} Jede Mini-Funktion ist ein eigener Service. Kopplung wird unbeherrschbar.
  \item \fachbegriff{Unklare Verträge.} Input/Output sind unscharf dokumentiert -- jeder Aufrufer rät.
  \item \fachbegriff{Kein Owner.} Der Service ist da, aber niemand pflegt ihn. Veraltet schnell.
\end{enumerate}

% --- 5. ORCHESTRIERUNG vs. CHOREOGRAPHIE ---------------------------------
\section{Orchestrierung vs.\ Choreographie}

Zwei Muster regeln, wie Services zusammenarbeiten:

\begin{figure}[h]
\centering
\begin{tikzpicture}[font=\sffamily\footnotesize,
  svc/.style={rectangle, rounded corners=2pt, draw=lpGray, fill=lpGrayLight,
    line width=0.5pt, minimum width=11mm, minimum height=7mm, inner sep=2pt},
  dir/.style={rectangle, rounded corners=2pt, draw=lpAccent, fill=lpAccentLight,
    line width=0.6pt, minimum width=18mm, minimum height=8mm, align=center, inner sep=2pt},
  arr/.style={-{Latex[length=1.8mm]}, draw=lpAccent, line width=0.5pt},
  ev/.style={-{Latex[length=1.8mm]}, draw=lpGray, line width=0.5pt, dashed},
  lbl/.style={font=\sffamily\scriptsize\bfseries, color=lpAccent}]
  % Orchestrierung
  \node[lbl] at (0,1.9) {Orchestrierung};
  \node[dir] (d) at (0,0.8) {Dirigent\\\scriptsize (Engine)};
  \node[svc] (a) at (-1.7,-0.6) {A};
  \node[svc] (b) at (0,-0.6) {B};
  \node[svc] (c) at (1.7,-0.6) {C};
  \draw[arr] (d) -- (a);
  \draw[arr] (d) -- (b);
  \draw[arr] (d) -- (c);
  % Choreographie
  \node[lbl] at (6,1.9) {Choreographie};
  \node[svc] (e) at (4.6,0.8) {A};
  \node[svc] (f) at (7.4,0.8) {B};
  \node[svc] (g) at (4.6,-0.8) {C};
  \node[svc] (h) at (7.4,-0.8) {D};
  \draw[ev] (e) -- (f);
  \draw[ev] (f) -- (h);
  \draw[ev] (h) -- (g);
  \draw[ev] (g) -- (e);
  \draw[ev] (e) -- (h);
\end{tikzpicture}
\caption{\sffamily Zentrale Steuerung (links) vs.\ ereignisbasierte Kopplung (rechts).}
\label{fig:konzept-71t03}
\end{figure}

\subsection{Orchestrierung}

In der \fachbegriff{Orchestrierung} steuert eine zentrale Komponente (der ``Dirigent'') den Ablauf -- typischerweise eine Workflow-Engine oder ein BPMN-Modell. Sie ruft die Services in der richtigen Reihenfolge auf, reagiert auf Antworten und steuert Verzweigungen.

\textbf{Vorteile:} klare Kontrolle, einfaches Monitoring (alles geht über den Dirigenten), gute Audit-Trails. \\
\textbf{Nachteile:} Single Point of Failure (wenn der Dirigent ausfällt), kann zum Engpass werden.% Korr: Grammatik "kann zur Engpass" -> "kann zum Engpass" (Maskulinum)

\subsection{Choreographie}

In der \fachbegriff{Choreographie} kommunizieren die Services direkt untereinander -- typischerweise über Events. Es gibt keinen zentralen Dirigenten; jeder Service ``weiß'', was er zu tun hat, wenn ein bestimmtes Ereignis eintritt.

\textbf{Vorteile:} skaliert sehr gut, keine zentrale Engpass-Komponente, flexibel erweiterbar. \\
\textbf{Nachteile:} schwer zu überblicken (kein zentrales ``Big Picture''), Debugging und Audit komplex.

\subsection{Wann was?}

\begin{tabularx}{\linewidth}{@{}>{\bfseries}lX@{}}
\toprule
Orchestrierung passt & Choreographie passt \\
\midrule
Klar definierte Geschäftsprozesse mit wenigen Varianten     & Lose gekoppelte Ökosysteme, viele unabhängige Reaktionen \\
Hohe Audit-/Compliance-Anforderungen                         & Sehr hohe Skalierungs-Anforderungen \\
Übersichtliche Anzahl Services (\textless\,30)               & Sehr viele Services oder externe Beteiligte \\
Steuerung durch Fachbereich (BPMN als Modell)                & Reine IT-getriebene Event-Architekturen \\
\bottomrule
\end{tabularx}

\medskip

\noindent In den meisten B2B-Unternehmen ist \emph{Orchestrierung} der pragmatischste Einstieg. Choreographie wird sinnvoll, wenn Skalierung und Entkopplung zu Problemen werden, die Orchestrierung nicht lösen kann.

% --- 6. WORKFLOW MANAGEMENT ----------------------------------------------
\section{Workflow Management Systeme als Steuerungsinstrument}

Ein \fachbegriff{Workflow Management System (WMS)} ist die Software, die Tasks aktiv steuert -- nicht nur dokumentiert, sondern \emph{ausführt}. Die wissenschaftliche Grundlage geht auf van der Aalst und van Hee zurück. \quelle{van der Aalst \& van Hee, 2002}

\subsection{Die sechs Workflow-Bausteine}

\begin{enumerate}
  \item \fachbegriff{Aufgaben (Tasks)} -- die einzelnen Schritte, die auszuführen sind.
  \item \fachbegriff{Regeln} -- wann welche Aufgabe ausgelöst wird, in welcher Reihenfolge, mit welchen Bedingungen.
  \item \fachbegriff{Rollen} -- wer welche Aufgabe ausführen darf oder soll.
  \item \fachbegriff{Eskalationen} -- was passiert, wenn eine Aufgabe nicht rechtzeitig erledigt wird (Trigger, Empfänger, Maßnahme).
  \item \fachbegriff{SLAs (Service Level Agreements)} -- definierte Service-Levels: Antwortzeit, Bearbeitungszeit, Qualität.
  \item \fachbegriff{Audit Trail / Monitoring} -- nachvollziehbare Aufzeichnung jeder Aktion, mit Dashboards zur aktuellen Sicht.
\end{enumerate}

Diese sechs Bausteine zusammen machen einen Prozess steuerbar. Fehlt einer, hat der Workflow Lücken: ohne Eskalationen versanden Aufgaben, ohne SLAs gibt es keine Erwartungs-Klarheit, ohne Audit Trail keine Nachweisbarkeit.

\subsection{Welchen Nutzen ein WMS bietet}

\begin{itemize}
  \item \textbf{Transparenz} -- jeder im Prozess sieht, wo welcher Vorgang gerade steht.
  \item \textbf{Durchlaufzeit-Steuerung} -- Wartezeiten werden sichtbar und adressierbar.
  \item \textbf{Qualitätssicherung} -- Gates und Checks sind erzwungen, nicht freiwillig.
  \item \textbf{Nachweisbarkeit} -- für Audits, Compliance und Reviews liegt eine vollständige Spur vor.
  \item \textbf{Priorisierung} -- bei knappen Ressourcen werden Aufgaben nach klaren Regeln zugewiesen.
\end{itemize}

\subsection{Wann ein WMS nicht hilft}

Ein WMS ist nicht für jeden Prozess sinnvoll:

\begin{itemize}
  \item \fachbegriff{Hohe Varianz.} Wenn fast jeder Vorgang anders ist, wird der Workflow zur Sammlung von Sonderfällen.
  \item \fachbegriff{Niedrige Frequenz.} Ein Prozess, der dreimal im Jahr läuft, ist ein Kandidat für eine Checkliste, nicht für einen Workflow.
  \item \fachbegriff{Hohe kreative Komponente.} ``Konzept entwickeln'' lässt sich nicht in feste Tasks zerlegen.
\end{itemize}

\begin{merkbox}[Workflow-Disziplin]
Ein Workflow ohne klare SLAs ist eine Wunschliste. Ein Workflow ohne Eskalationen ist eine Sammlung optionaler Aufgaben. Ein Workflow ohne Audit Trail ist nicht steuerbar -- nur beobachtbar. Wer ein WMS einführt, muss alle drei Elemente von Anfang an mitdenken.
\end{merkbox}

% --- 7. FALLSTUDIE -------------------------------------------------------
\section{Fallstudie: AeroTech Components}

Wir wenden den Tag auf ein konkretes Szenario an. \emph{AeroTech Components} entwickelt Komponenten für Industriekunden. Das Problem: Viele Ideen, chaotische Priorisierung, Prototypen verzögern sich, Qualitätsnachweise fehlen. Kunden drohen mit Lieferantenwechsel.

\subsection{Rahmenbedingungen}

\begin{itemize}
  \item Zeit: 10 Wochen für einen stabilen Idee-zu-Freigabe-Prozess.
  \item Budget: 120{.}000\,\euro{} (Workflow-Tool möglich, aber IT-Kapazität unklar).
  \item Zielkonflikte: Speed vs.\ Compliance, Standardisierung vs.\ Innovationsfreiheit.
  \item Datenlage: Zeitstempel aus Ticketing vorhanden, Qualitätsdaten verteilt.
\end{itemize}

\subsection{Schritt 1 -- Stakeholder-Analyse}

Top-8 Stakeholder + 2 externe:

\begin{tabularx}{\linewidth}{@{}>{\bfseries}lXll@{}}
\toprule
Stakeholder & Interesse & Einfluss & Strategie \\
\midrule
F\&E-Leitung           & Innovationsfreiheit erhalten            & hoch   & Mitentscheiden (Gate-Design) \\
Produktmanagement      & Business Case und Priorisierung         & hoch   & Mitentscheiden \\
Qualitätsmanagement    & Compliance und Nachweisbarkeit          & hoch   & Mitentscheiden \\
Operations/Produktion  & Reproduzierbarkeit                       & mittel & Konsultieren \\
Vertrieb / Key Account & Speed und Kundennutzen                  & hoch   & Konsultieren \\
IT / EA                & Architektur, Integrationen, Kapazität   & hoch   & Mitentscheiden \\
Einkauf                & Lieferanten und Kosten                  & mittel & Informieren \\
Finance / Controlling  & Wirtschaftlichkeit                      & mittel & Konsultieren \\
\midrule
Schlüsselkunde (extern) & Liefersicherheit, Qualität              & hoch   & Informieren \\
Auditor (extern)       & Nachweisbarkeit, Compliance              & hoch   & Informieren \\
\bottomrule
\end{tabularx}

\medskip

\noindent \textbf{Stopper-Stakeholder:} IT (wegen Kapazität/Integrationen) und F\&E-Leitung (wegen Innovationsfreiheit). Beide müssen früh und aktiv eingebunden werden -- nicht informiert, sondern mitentscheidend.

\subsection{Schritt 2 -- KPI-Framework}

Sechs KPIs über die vier Ebenen verteilt:

\begin{tabularx}{\linewidth}{@{}>{\bfseries}lXl@{}}
\toprule
KPI & Was misst sie & Ebene \\
\midrule
Lead Time Idee \textrightarrow{} Freigabe & Median in Wochen & Prozess (Lagging) \\
Rework-Rate Prototyp                    & \% der Prototypen mit Wiederholung & Quality (Leading/Lagging) \\
Gate-Pass-Rate beim ersten Durchlauf    & \% der Cases ohne Rückläufer im Gate & Quality (Leading) \\
Anteil ``Ideen mit Business Case''       & \% der eingereichten Ideen & Steuerungsqualität (Leading) \\
Kundenfeedback-Score nach Pilot          & Skala 1--10                & Outcome (Lagging) \\
Kosten pro freigegebenem Konzept         & \euro{} / Konzept           & Wirtschaftlichkeit (Lagging) \\
\bottomrule
\end{tabularx}

\medskip

\noindent \textbf{KPI-Steckbrief Beispiel ``Lead Time'':} Definition = Median über alle freigegebenen Cases im Zeitraum; Formel = Median(Timestamp Gate-Freigabe -- Timestamp Idee-Eingang); Quelle: Ticketing; Frequenz: wöchentlich; Owner: PMO; Ziel: $-20\,\%$ in 10 Wochen; Manipulationsrisiko: ``Schnell schließen ohne Qualität''; Gegenmaßnahme: Kopplung mit Gate-Pass-Rate und Rework-Rate. Dadurch ist die KPI-Beschleunigung nur dann ein Erfolg, wenn auch Qualitäts-Kennzahlen halten.

\subsection{Schritt 3 -- SOA-Service-Kandidaten}

Aus dem Prozess lassen sich fünf Service-Kandidaten ableiten:

\begin{tabularx}{\linewidth}{@{}>{\bfseries}lXll@{}}
\toprule
Service & Input \textrightarrow{} Output & Owner & Wiederverwendung \\
\midrule
Idee erfassen            & Idee, Kunde, Nutzen \textrightarrow{} Idea-ID                & PM        & in mehreren Prozessen \\
Business Case bewerten   & Nutzen, Kosten, Risiko \textrightarrow{} Score / Entscheidung & Finance/PM & überall wo BC nötig \\
Prototyp planen          & Anforderungen \textrightarrow{} Plan / SLA                   & F\&E/Ops  & spezifisch \\
Qualitätsnachweis erzeugen & Tests \textrightarrow{} Evidence-Paket                     & QM        & in mehreren Prozessen \\
Freigabe orchestrieren   & Evidence + Score \textrightarrow{} Freigabe / Reject         & Board     & Governance-übergreifend \\
\bottomrule
\end{tabularx}

\subsection{Schritt 4 -- Workflow-Skizze}

Soll-Logik (textuell):

\begin{quote}
\textbf{Start} ``Idee eingereicht'' \textrightarrow{} \emph{Task PM: Vollständigkeit prüfen} (\textbf{SLA 48\,h}) \\
\textrightarrow{} \textbf{XOR} ``Vollständig?'' \\
\qquad Nein \textrightarrow{} Rückfrage an Einreicher (\textbf{SLA 48\,h}, dann \textbf{Eskalation}) \\
\qquad Ja \textrightarrow{} \emph{Task Finance: Business Case Score} \\
\textrightarrow{} \textbf{XOR} ``Score $\geq$ Schwelle?'' \\
\qquad Nein \textrightarrow{} Backlog, Ende \\
\qquad Ja \textrightarrow{} \emph{Task F\&E: Prototyp planen} (\textbf{SLA 5 Tage}) \\
\textrightarrow{} \emph{Task QM: Testplan definieren} \\
\textrightarrow{} \emph{Task F\&E: Prototyp bauen} \\
\textrightarrow{} \emph{Task QM: Test durchführen} \\
\textrightarrow{} \textbf{XOR} ``Test bestanden?'' \\
\qquad Nein \textrightarrow{} Rework + \textbf{Eskalation nach 2 Schleifen} \\
\qquad Ja \textrightarrow{} \emph{Task Board: Freigabe} \\
\textrightarrow{} \textbf{Ende} ``Freigegeben''
\end{quote}

\subsection{Schritt 5 -- Maßnahmen unter Restriktionen (10 Wochen, IT-Unsicherheit)}

\begin{enumerate}
  \item \fachbegriff{Minimum-Workflow im bestehenden Ticketing} -- Regeln, SLAs und Eskalationen werden in der vorhandenen Ticketing-Software umgesetzt. Kein neues Tool, kein IT-Großprojekt.
  \item \fachbegriff{KPI-Minimum mit zwei Pflichtfeldern} -- jedes Ticket bekommt zwei zusätzliche Pflichtfelder (z.\,B.\ ``Business-Case-Status'' und ``Gate-Stufe''). Damit werden die KPIs Lead Time, Pass-Rate und Business-Case-Anteil messbar -- ohne neue Daten­integration.
\end{enumerate}

\begin{merkbox}[Wirkmechanismus der zwei Maßnahmen]
Die erste Maßnahme greift \emph{strukturell}: sie etabliert SLAs und Eskalationen, die das Versanden verhindern. Die zweite Maßnahme greift \emph{datenbasiert}: zwei zusätzliche Felder machen die Wirkung messbar. Zusammen schaffen sie Steuerbarkeit \emph{ohne} Tool-Großprojekt -- die robusteste Wirkung unter den Restriktionen.
\end{merkbox}

% --- 8. COACHING ---------------------------------------------------------
\section{Coaching: KPI-Hygiene und Verantwortungsfähigkeit}

Coaching auf Senior-Niveau führt von ``KPIs einführen'' zu ``KPIs hygienisch betreiben'' -- mit klarer Verantwortlichkeit und aktivem Schutz vor Fehlsteuerung.

\subsection{Sechs Reflexionsfragen vor jeder KPI-Einführung}

\begin{enumerate}
  \item Welche Entscheidung soll die KPI stützen -- und was würden wir konkret anders tun, wenn sich der Wert verdoppelt oder halbiert?
  \item Welches Verhalten wird die KPI auslösen -- gewünscht oder unerwünscht?
  \item Wer ist Owner der KPI -- und ist der Owner zugleich derjenige, der das gemessene Verhalten beeinflussen kann?
  \item Welche Manipulationsmöglichkeit gibt es -- und durch welche Gegen-KPI schützen wir uns?
  \item Wie lange wird die KPI gültig sein -- gibt es einen Reviewzyklus, in dem sie überprüft wird?
  \item Welche Stakeholder sind betroffen -- und wie sichern wir, dass sie die KPI als fair und nützlich akzeptieren?
\end{enumerate}

\subsection{Konfliktkompetenz statt Konflikt-Vermeidung}

Senior-Modellierer verdrängen Stakeholder-Konflikte nicht -- sie machen sie sichtbar. Die Disziplin:

\begin{itemize}
  \item \fachbegriff{Konflikt benennen} -- ``Wir haben hier einen Trade-off zwischen Speed und Compliance.''
  \item \fachbegriff{Beide Seiten würdigen} -- ``Beide Anforderungen sind legitim. Vollkompromiss gibt es nicht.''
  \item \fachbegriff{Entscheidung herbeiführen} -- ``Wir entscheiden uns für Speed in der Pilotphase, Compliance wird mit zwei Prüf-Gates abgesichert.''
  \item \fachbegriff{Verantwortung klären} -- ``Owner dieser Entscheidung ist X. Falls die Trade-offs sich verschieben, eskaliert X an Y.''
\end{itemize}

Wer Konflikte nur ``wegmoderiert'' (in Form von ``wir machen beides''), erzeugt unklare Verantwortlichkeiten -- und damit später Streit unter veränderten Bedingungen.

\subsection{Senior-Shift: vom Modellieren zum Steuern}

Drei Kennzeichen, an denen man den Wechsel erkennt:

\begin{itemize}
  \item Man fragt \emph{vor} der Modellierung: ``Welche Entscheidung soll dieses Modell stützen?'' -- statt ``welche Notation nehmen wir?''
  \item Man ist bereit, KPIs zu \emph{ändern}, wenn sie zur Fehlsteuerung führen -- statt sie zu verteidigen.
  \item Man führt Diskussionen mit dem Vorstand über die \emph{Steuerungslogik} -- nicht über die Notation.
\end{itemize}

\begin{eselsbox}[Eselsbrücke: \akzent{Z-V-M-G}]
Vier Senior-Hebel in der Prozess-Steuerung:\\[2pt]
\textbf{Z}weck -- welche Entscheidung wird gestützt?\\
\textbf{V}erantwortung -- wer trägt die Entscheidung?\\
\textbf{M}anipulationsrisiko -- wie schützen wir uns gegen Fehlsteuerung?\\
\textbf{G}egen-KPI -- welche Kopplung sichert die Hauptmetrik ab?\\[2pt]
\emph{``Z-V-M-G'' -- vier Fragen, die jede KPI-Einführung überlebensfähig machen.}
\end{eselsbox}

% --- 9. MANAGEMENT ------------------------------------------------------
\section{Management-Perspektive: Steuerungsarchitektur}

Auf Wissens- und Anwendungs-Ebene haben wir die vier Hebel (Stakeholder, KPIs, SOA, Workflow) einzeln betrachtet. Auf Management-Ebene geht es um die Frage: \emph{Wie verankern wir sie als zusammen­hängende Architektur?}

\subsection{Die Steuerungsarchitektur auf einer Seite}

Eine wirksame Steuerungsarchitektur beantwortet fünf Fragen auf einer Seite:

\begin{enumerate}
  \item \textbf{Welche Entscheidungen sollen unsere KPIs und Workflows tatsächlich stützen?} -- Priorisierung, Gate-Freigabe, Ressourcen-Allokation, Risiko-Eskalation, Compliance-Nachweis.
  \item \textbf{Welche KPI-Klassifikation gilt?} -- Outcome, Output, Prozess, Compliance/Risiko -- mit klaren Definitions-Standards.
  \item \textbf{Welche Rollen-Verantwortung gibt es?} -- RACI für KPI-Änderungen, Workflow-Änderungen, Service-Katalog.
  \item \textbf{Welche Service-Kriterien?} -- wann lohnt ein Service, welcher Vertrag (Input/Output), welche Integrationsprinzipien (API-First / Events)?
  \item \textbf{Welches Workflow-Betriebsmodell?} -- SLAs, Eskalationspfade, Audit Trail, Definition of Done für Workflow-Änderungen.
\end{enumerate}

Wenn diese fünf Punkte nicht auf einer Seite Platz finden, ist die Architektur zu kompliziert -- nicht zu umfangreich.

\subsection{Die KPI-Policy als Goodhart-Schutz}

Eine KPI-Policy regelt die Spielregeln, unter denen KPIs eingeführt, geändert und abgeschaltet werden:

\begin{itemize}
  \item Jede neue KPI braucht einen vollständigen Steckbrief mit Manipulations\-risiko-Analyse.
  \item Jede KPI hat eine vorab definierte Gegen-KPI oder Kopplung.
  \item Jede KPI wird mindestens jährlich auf Wirksamkeit überprüft.
  \item Wer eine KPI einführen will, muss vorab die Frage beantworten: ``Welche Entscheidung wird sich dadurch konkret ändern?''
\end{itemize}

Ohne KPI-Policy entstehen über Jahre Hunderte Kennzahlen -- die wenigen wirksamen ertrinken im Reporting-Rauschen.

\subsection{Drei Executive-KPIs als Auswahl}

Bei der Vielzahl möglicher KPIs muss das Management drei bis fünf ``Executive KPIs'' definieren, die im Vorstand oder Aufsichtsrat regelmäßig diskutiert werden. Auswahl-Kriterien:

\begin{itemize}
  \item Trade-off-sichtbar machend -- die KPI zeigt einen tatsächlichen Trade-off (z.\,B.\ Lead Time und Rework gekoppelt).
  \item Strategisch relevant -- die KPI beeinflusst die wichtigsten Geschäftsziele.
  \item Manipulationssicher -- es gibt eine Gegen-KPI, die die Hauptmetrik schützt.
  \item Stabil -- die KPI bleibt über mehrere Quartale relevant, ohne ständig zu ändern.
\end{itemize}

\subsection{Anti-Patterns auf Management-Ebene}

\begin{enumerate}
  \item \fachbegriff{``Wir messen alles -- mehr Daten, bessere Entscheidungen.''} Falsch: mehr Kennzahlen erzeugen mehr Aufwand, nicht mehr Erkenntnis.
  \item \fachbegriff{``KPIs sind Sache des Controllings.''} Falsch: KPIs sind Sache der Prozess-Owner. Controlling ist Werkzeug, nicht Inhaber.
  \item \fachbegriff{``Wir nennen unsere Services ‚Microservices'.''} Falsch: Naming ändert nichts. Wichtig sind Kohäsion und Kopplung.
  \item \fachbegriff{``Der Workflow wird im Tool definiert.''} Falsch: erst die Logik (BPMN als Soll), dann das Tool. Wer im Tool startet, baut Tool-Speziallösungen statt Geschäftsprozesse.
\end{enumerate}

\begin{merkbox}[Schluss-Merksatz für Tag 3]
Steuerung ist eine \textbf{Architektur­entscheidung} -- keine Tool-Entscheidung. Sie integriert Stakeholder-Klarheit, KPI-Hygiene, Service-Disziplin und Workflow-Governance zu einem Operating Model. Wenn alle vier Hebel sauber zusammenspielen, wird aus einem Diagramm ein lebender Steuerungs­apparat. Fehlt einer, bleibt es Papier.
\end{merkbox}

% --- 10. HAEUFIGE FEHLER ------------------------------------------------
\section{Häufige Fehler in der Prozess-Steuerung}

\subsection{Sechs typische Anti-Muster}

\begin{enumerate}
  \item \fachbegriff{Stakeholder als Pflichttermin behandeln.} ``Wir haben sie informiert.'' -- aber nicht eingebunden. Die Stopper-Stakeholder erfahren zu spät und blockieren spät.
  \item \fachbegriff{KPI ohne Steckbrief einführen.} Die Zahl wird gemessen, aber niemand weiß, wer was tun soll, wenn sie sich ändert.
  \item \fachbegriff{Goodhart ignorieren.} Eine Hauptmetrik ohne Gegen-KPI führt systematisch zur Verzerrung.
  \item \fachbegriff{Servicegrenzen technisch ziehen, nicht fachlich.} Services entstehen entlang IT-Systeme statt entlang fachlicher Fähigkeiten -- Kohäsion sinkt, Kopplung steigt.
  \item \fachbegriff{Workflow ohne SLAs.} Aufgaben werden im Tool dokumentiert, aber niemand muss sie in einer bestimmten Zeit erledigen. Ergebnis: alles dauert ``so lange wie es eben dauert''.
  \item \fachbegriff{Audit Trail zu spät einführen.} Wenn der Audit Trail erst kommt, nachdem ein Vorfall passiert ist, ist die Spur weg. Audit Trail gehört von Tag 1 ins System.
\end{enumerate}

\subsection{Review-Checkliste vor einer Workflow-Einführung}

\begin{enumerate}
  \item Sind die Top-5 Stakeholder benannt und mit klarer Einbindungsstufe versehen?
  \item Ist mindestens eine ``Stopper-Stakeholder''-Strategie definiert?
  \item Hat jede zentrale KPI einen vollständigen Steckbrief mit Gegen-KPI?
  \item Sind alle Services mit Input/Output/Owner dokumentiert?
  \item Hat der Workflow definierte SLAs, Eskalationen und einen Audit Trail von Tag 1 an?
\end{enumerate}

% --- 11. TAKE-AWAYS ------------------------------------------------------
\section{Take-aways aus Tag 3}

\begin{enumerate}
  \item \textbf{Vier Hebel der Steuerung.} Stakeholder, KPIs, SOA-Services und Workflow-Orchestrierung -- alle vier müssen zusammenpassen.
  \item \textbf{Stakeholder-Matrix.} Einfluss $\times$ Interesse -- vier Quadranten, vier Strategien. Stopper-Stakeholder zuerst identifizieren.
  \item \textbf{KPI-Architektur in vier Ebenen.} Outcome, Output, Prozess, Compliance/Risiko -- mit Leading und Lagging Indicators ausgewogen.
  \item \textbf{Goodhart-Schutz.} Jede zentrale KPI braucht eine Gegen-KPI. Sonst entsteht Fehlsteuerung.
  \item \textbf{KPI-Steckbrief in 10 Feldern.} Ohne vollständigen Steckbrief ist es keine KPI.
  \item \textbf{Servicegrenzen mit Kohäsion und Kopplung.} Hohe Kohäsion, geringe Kopplung -- das 50-Jahre-Prinzip gilt unverändert.
  \item \textbf{Orchestrierung vs.\ Choreographie.} Im B2B-Umfeld meist Orchestrierung als pragmatischer Einstieg.
  \item \textbf{Workflow mit sechs Bausteinen.} Aufgaben, Regeln, Rollen, Eskalationen, SLAs, Audit Trail -- alle sechs müssen vorhanden sein.
  \item \textbf{Steuerung ist Architektur, nicht Tool.} Erst die Logik, dann das System.
\end{enumerate}

% --- 12. LITERATUR -------------------------------------------------------
\clearpage
\section{Literatur}

Im Skript verwendete und für die Vertiefung empfohlene Quellen. Zitierstil: APA-7.

\begin{description}[style=nextline,leftmargin=18pt,labelindent=0pt,itemsep=4pt]
  \item[Dumas, M., La Rosa, M., Mendling, J., \& Reijers, H.\,A. (2018).] \emph{Fundamentals of business process management} (2nd ed.). Springer. \href{https://doi.org/10.1007/978-3-662-56509-4}{https://doi.org/10.1007/978-3-662-56509-4}
  \item[Erl, T. (2005).] \emph{Service-oriented architecture: Concepts, technology, and design}. Prentice Hall.
  \item[Freeman, R.\,E. (1984).] \emph{Strategic management: A stakeholder approach}. Pitman.
  \item[Kaplan, R.\,S., \& Norton, D.\,P. (1996).] \emph{The balanced scorecard: Translating strategy into action}. Harvard Business School Press.
  \item[Mendelow, A.\,L. (1981).] Environmental scanning -- The impact of the stakeholder concept. \emph{Proceedings of the International Conference on Information Systems (ICIS 1981)}, Paper 20.
  \item[Parmenter, D. (2015).] \emph{Key performance indicators: Developing, implementing, and using winning KPIs} (3rd ed.). Wiley. \href{https://doi.org/10.1002/9781119019855}{https://doi.org/10.1002/9781119019855}
  \item[Project Management Institute. (2021).] \emph{A guide to the Project Management Body of Knowledge (PMBOK\textregistered{} Guide)} (7th ed.). Project Management Institute.
  \item[Strathern, M. (1997).] ``Improving ratings'': Audit in the British university system. \emph{European Review, 5}(3), 305--321. (häufig zitierte Formulierung des Goodhart-Effekts)
  \item[van der Aalst, W.\,M.\,P., \& van Hee, K.\,M. (2002).] \emph{Workflow management: Models, methods, and systems}. MIT Press.
\end{description}

% --- 13. GLOSSAR ---------------------------------------------------------
\clearpage
\section{Glossar}

Die wichtigsten Begriffe des Tages -- kurz definiert, in alphabetischer Reihenfolge.

\begin{description}[style=nextline,leftmargin=0pt,labelindent=0pt,itemsep=4pt]
  \item[Audit Trail] Lückenlose, nachvollziehbare Aufzeichnung aller Aktionen in einem Workflow. Voraussetzung für Compliance und Reviews.
  \item[Choreographie] Architekturmuster, bei dem Services über Events direkt miteinander kommunizieren -- ohne zentralen Dirigenten.
  \item[Compliance-KPI] Kennzahl, die regulatorische, vertragliche oder interne Vorgaben misst (z. B. Audit-Findings, Fehlerquote).
  \item[Eskalation] Vorab definierter Trigger plus Maßnahme, wenn eine Aufgabe nicht termingerecht erledigt wird.
  \item[Goodhart-Effekt] Beobachtung, dass eine Kennzahl ihre Aussagekraft verliert, sobald sie zum Ziel wird. \quelle{nach Goodhart, in Strathern, 1997}
  \item[Kohäsion] Architekturprinzip: alles in einem Service gehört fachlich zusammen, der Service tut eine Sache.
  \item[Kopplung] Architekturprinzip: Services sollen möglichst unabhängig sein, Änderungen in einem brechen nicht zwangsläufig andere.
  \item[KPI-Steckbrief] Vollständige Definition einer Kennzahl in 10 Feldern (Name, Definition, Formel, Quelle, Frequenz, Owner, Ziel, Verhalten, Manipulation, Gegenmaßnahme). \quelle{Parmenter, 2015}
  \item[Lagging Indicator] Kennzahl, die ein Ergebnis misst (vergangenes Verhalten). Präzise, aber spät.
  \item[Leading Indicator] Kennzahl, die einen Vorboten misst (zukünftiges Verhalten). Ungenauer, aber rechtzeitig.
  \item[Mendelow-Matrix] 2$\times$2-Matrix zur Stakeholder-Klassifikation (Einfluss $\times$ Interesse). \quelle{Mendelow, 1981}
  \item[Orchestrierung] Architekturmuster, bei dem eine zentrale Komponente (Engine) die Service-Aufrufe steuert.
  \item[Outcome-KPI] Kennzahl, die den Kundennutzen misst (z. B. NPS, Kundenzufriedenheit).
  \item[Output-KPI] Kennzahl, die die produzierte Menge misst (z. B. Anzahl Freigaben).
  \item[Prozess-KPI] Kennzahl, die die Effizienz des Ablaufs misst (z. B. Lead Time, Rework).
  \item[SLA -- Service Level Agreement] Vereinbarung über Service-Levels: Antwortzeit, Bearbeitungszeit, Qualität.
  \item[Service (SOA)] Klar abgegrenzte Fähigkeit mit definiertem Vertrag (Input/Output) und Owner. \quelle{Erl, 2005}
  \item[Stakeholder] Person oder Gruppe, die ein Vorhaben beeinflussen kann oder davon betroffen ist. \quelle{Freeman, 1984}
  \item[Stopper-Stakeholder] Stakeholder mit hohem Einfluss, deren Ablehnung das Vorhaben stoppen kann -- müssen aktiv eingebunden werden.
  \item[Workflow Management System (WMS)] Software, die Tasks aktiv steuert (Aufgaben, Regeln, Rollen, Eskalationen, SLAs, Audit Trail). \quelle{van der Aalst \& van Hee, 2002}
\end{description}

\vspace{1cm}
\noindent{\color{lpAccent}\rule{\linewidth}{0.6pt}}\\[4pt]
{\footnotesize\sffamily\color{lpGray}Lernplatt \textperiodcentered{} by Can Siebert \textperiodcentered{} Kurs 71 (Dig 42) \textperiodcentered{} Tag 3 \textperiodcentered{} \textcopyright{} 2026}

\end{document}
